\title{LongRoPE: Extending LLM Context Window Beyond 2 Million Tokens}

\begin{abstract}
A substantial context window is highly valued in large language models (LLMs). Nonetheless, present methods are constrained to approximately 128k tokens, primarily because of the significant computational expense of fine-tuning, the limited availability of lengthy textual data, and the disruptive effects that arise from newly introduced token positions.

In this work, we present LongRoPE, a novel approach that, for the first time, expands the context window of pre-trained LLMs to a remarkable \textbf{2048k} tokens. This is accomplished using as few as 1,000 fine-tuning steps within a training length of up to 256k, all while preserving the original performance at short context lengths. The success of LongRoPE relies on three principal contributions: (i) we recognize and leverage two distinct types of non-uniformity in positional interpolation via a computationally efficient search, resulting in superior initialization for fine-tuning and facilitating an eightfold extension in scenarios where fine-tuning is not applied; (ii) we propose a staged extension process, beginning with fine-tuning an LLM at 256k length, followed by a subsequent positional interpolation on this already-extended model to reach the full 2048k context length; (iii) we perform readjustment of LongRoPE at the 8k context window, thereby restoring performance for short-range tasks. Comprehensive evaluations on LLaMA2 and Mistral over diverse benchmarks validate the efficacy of our methodology. LongRoPE-expanded models largely preserve the original model architecture, introducing only minor adjustments to the positional encoding, and remain compatible with most established optimization techniques. The code will be released at \texttt{https://github.com/microsoft/LongRoPE}

\section{Introduction}

Although Large Language Models (LLMs) have achieved impressive performance across a range of tasks~\cite{gpt4,llama2}, they remain constrained by their context window size; for example, LLaMA2 is limited to 4096 tokens~\cite{llama2}. When required to process inputs that exceed this window, their effectiveness diminishes since these models are not exposed to such extended positions during training. This limitation is particularly problematic for key applications such as in-context learning with a large set of examples~\cite{Huang2023FewerIM} and the operation of LLM agents~\cite{park2023generative,madaan2023selfrefine}.

Recent research demonstrates that it is possible to expand the context window of a pre-trained LLM to approximately 128k through fine-tuning on longer input sequences~\cite{longlora,pi,yarn,zhang2024soaring,liu2023scaling}. Nonetheless, pushing the context window beyond this range faces three key challenges. Firstly, the incorporation of previously unused position indices tends to produce problematic values, which exacerbates out-of-distribution complications and significantly hampers the convergence of fine-tuning~\cite{pi}. This issue is heightened during substantial extensions, such as moving from 4k to more than $>$1000k, where over 90\% of the positional indices are newly introduced. Secondly, effective fine-tuning generally necessitates training samples that match the targeted context lengths. The availability of very long texts (particularly those surpassing 1000k) in current corpora is minimal. Additionally, training on such lengthy sequences imposes considerable computational demands, as it requires an extraordinary amount of both training time and GPU resources. Thirdly, when scaling to context windows of extreme length, the distribution of attention across a vast number of token positions becomes diluted, resulting in diminished performance on shorter contexts commonly encountered during inference~\cite{pi}.

A common method to address the initial challenge involves applying RoPE positional embedding interpolation~\cite{rope,pi}, which maps new position indices to the original pretrained interval, as illustrated in Fig.\ref{fig:overview}. Position Interpolation (PI)~\cite{pi} executes a linear interpolation of RoPE’s rotary phases based on the extension factor. The NTK method~\cite{ntk,dynamicntk} instead proposes non-uniform interpolation and extrapolation strategies along RoPE dimensions. In contrast, YaRN~\cite{yarn} organizes the RoPE dimensions into three distinct groups determined by their frequency characteristics, employing extrapolation, NTK-based, and linear interpolation techniques for each group, respectively. Nevertheless, within Transformer models, the positional embedding manifests \emph{highly non-uniform information entropy}. Existing methods do not adequately harness this fine-grained non-uniformity, resulting in a loss of information and consequently restricting the extent of the context window.

Section~\ref{sec:analysis} demonstrates two principal empirical observations: \textbf{(1)} Robust positional interpolation necessitates addressing two types of non-uniformity—specifically, the disparity among RoPE dimensions and the heterogeneity in token positions. Both the lower RoPE dimensions and tokens appearing at the start benefit from applying less interpolation; however, the optimal strategy depends on the desired sequence extension length. 
\textbf{(2)} Incorporating these non-uniform features into positional interpolation facilitates the preservation of critical information from the original RoPE, particularly within important dimensions and early token positions. This approach reduces the information loss typically associated with positional interpolation, resulting in improved initial conditions for subsequent fine-tuning. Additionally, this method enables up to 8$\times$ context window extension even without further fine-tuning.

Building upon these results, we propose \textbf{LongRoPE}, a powerful approach designed to stretch the LLM context window to over 2 \emph{million} tokens. LongRoPE is constructed around three core advancements. Firstly, {LongRoPE} leverages multidimensional, non-uniform strategies for positional interpolation to their full advantage. Specifically, it determines optimal rescale factors for the rotation angles used in RoPE within each RoPE dimension according to token position. Given that the space of potential rescale factors increases exponentially with the desired window extension, {LongRoPE} incorporates an evolutionary search procedure equipped with two optimization strategies to enhance the efficiency of this process. An example of a resulting rescaled RoPE configuration is depicted in Fig.~\ref{fig:overview}.

Subsequently, {LongRoPE} employs an efficient and incremental extension approach that enables reaching a 2048k context window, all without necessitating direct fine-tuning on ultra-long documents, which are both rare and difficult to obtain. The process commences by identifying an appropriate 256k sequence length for the existing pre-trained LLM and performing fine-tuning at this length. Due to our non-uniform positional interpolation scheme permitting an 8$\times$ expansion even in the absence of further fine-tuning, we proceed with a secondary search for updated RoPE rescaling factors using the previously fine-tuned, length-extended LLM. Through this process, the method successfully expands the context window to 2048k for both LLaMA2 and Mistral~\cite{mistral}.

To address potential drops in performance on the original, shorter context lengths, {LongRoPE} further tunes the RoPE rescaling parameters on the extended LLM. In a process analogous to scaling the model up from 256k to 2048k, our approach involves scaling the 256k fine-tuned LLM down to 4k and 8k context windows, utilizing our search algorithm to minimize positional interpolation effects. For inference tasks involving sequence lengths shorter than 8k, we apply the rescale factors identified by our search process to update RoPE.

A comprehensive series of experiments conducted on multiple LLMs and a range of long-context tasks substantiates the efficacy of our approach. Our findings indicate that LongRoPE consistently preserves low perplexity across evaluation lengths spanning from 4k up to 2048k, surpasses 90\% accuracy in passkey retrieval, and achieves accuracy levels on par with standard benchmarks set within a 4096 context window. The LongRoPE technique is compatible with any LLM utilizing RoPE embeddings. We intend to make both our code and the LongRoPE-2048k models publicly available.

\section{Non-uniformity in Positional Interpolation}
\label{sec:analysis}

\subsection{Preliminary}

Transformer models require explicit positional information, often in the form of position embedding, to represent the order of input tokens. Our work focuses on the RoPE~\cite{rope} position embedding, which is widely used in recent  LLMs. For a token at position index $n$, its corresponding RoPE encoding can be simplified as follows:

\begin{equation}
		\fontsize{7.8}{7.8} \selectfont
	\label{eq:rope}
	\left[ cos(n\theta_0), sin(n\theta_0), cos(n\theta_1), \cdot\cdot\cdot, cos(n\theta_{d/2-1}), sin(n\theta_{d/2-1}) \right]   
\end{equation}

where $d$ denotes the embedding dimension, 
$n\theta_i$ indicates the rotary angle corresponding to the token at position $n$, 
and $\theta_i=\theta^{-2i/d}$ specifies the rotational frequency parameters. In the standard configuration of RoPE, the base value for $\theta$ is set to 10000. 

\textbf{\emph{Context window extension ratio $s$} and positional interpolation}. Here, $s$ is introduced as the ratio between the new, extended context length $L'$ and the initial context length $L$: $s=\frac{L'}{L}$.

To extend the context window from $L$ to $L'$, current positional interpolation methods suggest downscaling rotation frequencies $\theta_i$ based on extension ratio $s$. Let $\beta=\theta^{2/d}$, and $\lambda$ denote the actual rescale factor related to $s$, we   unify these positional interpolation methods as follows:

\begin{equation}	
	\fontsize{7.5}{7.5} \selectfont
	\label{eq:unified_rope}
	\begin{aligned}
		\left[cos\left(\frac{n}{\lambda(\beta)^0}\right), sin \left(\frac{n}{\lambda(\beta)^0}\right), cos\left(\frac{n}{\lambda(\beta)^1}\right), \cdot\cdot\cdot,  sin \left(\frac{n}{\lambda(\beta)^{d/2-1}}\right) \right]   
	\end{aligned}
\end{equation}

\textbf{Linear positional interpolation (PI)}. The PI method~\cite{pi} involves linearly rescaling position indices when the input sequence is extended, but only within the length for which the model was originally trained. Specifically, given a target extension ratio $s$, the method uniformly compresses the rotation angles associated with all positions by a factor of $\lambda = s$ across each RoPE dimension. While straightforward, this approach results in positional representations being densely packed, making it challenging for the model to differentiate between nearby tokens. Consequently, PI often exhibits suboptimal performance, especially at larger extension ratios.

\textbf{NTK-based interpolation and extrapolation}. The studies in ~\cite{ntk,dynamicntk} reinterpret RoPE through the lens of information representation and utilize the Neural Tangent Kernel (NTK) theory~\cite{jacot2018neural,tancik2020fourier}. To address the issue of positional overlap present in PI, they propose a strategy that spreads interpolation demands over the different dimensions of RoPE. Specifically, this approach reduces scaling for dimensions with higher frequencies and increases it for those with lower frequencies, thereby facilitating both interpolation and extrapolation of positional encodings, where $\lambda=s^{i}$. The enhanced dynamic NTK~\cite{dynamicntk} further refines this method by varying the extension ratio at each position according to the sequence’s current length. In contrast to PI, which typically requires fine-tuning, NTK-aware techniques are capable of expanding context windows even without further fine-tuning, although they generally support a maximum extension ratio of 4$\times$.

\textbf{YaRN}~\cite{yarn} presents a notable advancement in positional interpolation. This method segments RoPE dimensions according to their frequency into three distinct groups, each assigned a specialized interpolation approach. Specifically, high frequency dimensions utilize extrapolation ($\lambda$=1), low frequency dimensions are treated with standard linear interpolation (PI), and mid-range frequency dimensions are processed using the NTK method. The main innovation of YaRN is this categorization of RoPE dimensions. However, this grouping currently relies on manual, experience-driven procedures, potentially resulting in less than optimal outcomes when applied to novel LLMs.

\subsection{Study on Non-uniform Positional Interpolation}

Building upon insights from NTK and YaRN, we observe that their improvements stem largely from introducing non-linearity, particularly in treating various frequencies along the RoPE dimensions to enhance both interpolation and extrapolation capabilities. Despite this, most non-linear approaches implemented so far depend on heuristics crafted by humans. This leads us to ask: (1) Could positional interpolation methods be further optimized? (2) Do alternative forms of non-linearity exist that have yet to be investigated?

To investigate these issues, we employ evolution search (refer to Sec.~\ref{sec:method}) to identify more effective non-uniform positional interpolations for LLaMA2-7B. Perplexity serves as the optimization objective during the search, evaluated on 5 randomly selected instances from the PG19~\cite{pg19} validation dataset. Our experimental study uncovers several important observations.

\textbf{Finding 1}: \textit{RoPE dimensions exhibit substantial non-uniformities, which are not effectively handled by current positional interpolation methods.}

We conduct a search for the optimal value of $\lambda$ for each RoPE dimension as described in Eq.~\ref{eq:unified_rope}. The results presented in Table~\ref{tbl:exp1} compare the perplexity of LLaMA2-7B across various methods on the PG19 and Proof-pile~\cite{proof-pile} test sets, all evaluated without additional fine-tuning. Our optimized configuration achieves substantial perplexity reductions, indicating that existing interpolation techniques—both linear (PI) and non-uniform (Dynamic-NTK and YaRN)—do not reach optimality. Interestingly, YaRN performs worse than both PI and NTK on the PG19 dataset, due to its inability to fully utilize the intended extended context window in a model that has not been fine-tuned. As a specific illustration, YaRN’s perplexity notably increases beyond 7k tokens when the context size is set to 8k.

In our investigation, we observe that the rescaling factors $\lambda$ in Eq.~\ref{eq:unified_rope} exhibit variability, unlike the constant scale $s$ utilized in PI, the NTK formula, or the group-based computation of YaRN. Employing these heterogeneous scaling factors leads to marked enhancements in LLaMA2's language modeling ability—as measured by perplexity—across 8k and 16k context window lengths, all achieved without additional fine-tuning. The underlying reason is that the adapted positional embeddings remain closely aligned with the original RoPE structure, particularly safeguarding critical dimensions. This alignment facilitates more precise differentiation between proximate token positions, thereby reducing the cognitive load on the LLM.

\textbf{Finding 2}: \textit{RoPE for the initial tokens in the input sequence should be extrapolated with less interpolation.} 

We propose that for the first $\hat{n}$ tokens of the input, position interpolation via RoPE should be minimized. These tokens typically receive higher attention weights and play a more significant role in attention computations, a pattern noted in works such as Streaming LLM~\cite{streamingllm} and LM-Infinite~\cite{han2023lminfinite}. To evaluate this idea, we expand the context window to sizes of 8k and 16k with PI and NTK while systematically excluding the initial $\hat n$ (where $\hat n=0,2,\ldots,256$) tokens from interpolation. When $\hat n = 0$, the standard PI and NTK procedures are recovered. Table~\ref{tbl:tokenposition} presents two primary findings: \textbf{(1)} Excluding the leading tokens from position interpolation leads to measurable gains in performance; and \textbf{(2)} the best value of $\hat n$ varies according to the extended context window length.

\textbf{Finding 3}: \textit{Non-uniform positional interpolation effectively extends LLM context window in both fine-tuning and non-fine-tuning settings.} 

Although our non-uniform position interpolation method already yields considerable gains for context extensions up to 8k and 16k tokens without requiring additional training, adapting to even longer context lengths necessitates fine-tuning. Therefore, we perform fine-tuning on LLaMA2-7B using our optimized RoPE strategy to support a 64k token context window (details in Appendix). As illustrated in Table~\ref{tbl:finetune}, our approach achieves superior results compared to PI and YaRN, both prior to and following fine-tuning of LLaMA2-7B. This enhanced performance stems from more effective non-uniform positional interpolation, which not only reduces information loss but also yields a stronger initialization for the fine-tuning process.

\textbf{Summary}. Our findings reveal two types of non-uniformity: those associated with RoPE dimensions and with token positions. Leveraging these forms of non-uniformity within positional interpolation substantially advances LLM context length extension capabilities.

\section{LongRoPE}

\label{sec:method}
Building on these insights, we propose LongRoPE, a framework that initially employs an efficient search strategy designed to leverage both identified non-uniformities. This approach enables the extension of the LLM context window to surpass 2 million tokens.

\subsection{Problem Formulation}

These two sources of non-uniformity expand the solution space considerably and add layers of complexity to the optimization process. To manage this challenge, we recast the multidimensional non-uniform position interpolation optimization as a search problem.

For a LLM targeting a  context window size of $L'$ and lengthy input documents $\mathbf{X}$, where each $\mathbf{x}\in\mathbf{X}$ surpasses $L'$ in token length, we denote the original rotary angle of the $i^{th}$ dimension in RoPE embedding at token position $n$ as  $\frac{n}{\beta^i}$. The optimization problem is then formulated as follows:

\begin{equation}
\begin{gathered}
		\label{eq:problem}

\underset {\mathbf{x}\in\mathbf{X}; \, |\mathbf{x}|\ge L'}{ \operatorname {arg\,min} } \,  \mathcal{L}\, \left( \text{LLM} (\text{RoPE}, \mathbf{X})\right), \text{where}
\\
\scalebox{0.93}{
$\underset {\begin{subarray}{l} i=0,\ldots,\frac{d}{2}-1;\\ n\in[0,|\mathbf{x}|);\end{subarray}}{\text{RoPE}_{(n)}} = \left[\ldots, \cos\left(\mathbb{I}(\hat{\lambda}_{i}, \hat{n}) \cdot \frac{n}{\beta^i}\right), \sin\left(\mathbb{I}(\hat{\lambda}_{i}, \hat{n}) \cdot \frac{n}{\beta^i}\right), \ldots\right]$}
\\
\text{with} \quad \mathbb{I}(\hat{\lambda}_{i},\hat{n}) = \begin{cases}
1 & \text{if $n < \hat{n}$} \\
\frac{1}{\lambda_i} & \text{if $n \geq \hat{n}$}
\end{cases}
\end{gathered}
\end{equation}

In this context, we define a collection of rescale factors, $\mathbb{I}(\hat{\lambda}_{i},\hat{n})$, designed to accommodate two types of non-uniformities. Here, $\hat{\lambda_i}$ represents the non-uniformity across RoPE dimensions, while $\hat{n}$ corresponds to non-uniformity in token positions. The function $\mathbb{I}(\hat{\lambda}_{i},\hat{n})$ is specifically utilized to adjust the rotation angle associated with the $i^{th}$ RoPE dimension, employing $\hat\lambda_{i}$ as the adjustment factor and $\hat{n}$ as a threshold for token positions. When considering the first $\hat{n}$-$1$ token positions, the adjustment factor $\hat\lambda_{i}$ is not applied, maintaining the standard RoPE rotary angle $\frac{n}{\beta^i}$. For positions where $n\ge\hat{n}$, the rescale factor is incorporated into the computation.

Assuming a desired context window length of $L'$, our goal is to determine the best rescale factors ($\mathbb{I}(\hat{\lambda}_{0},\hat{n})$, $\mathbb{I}(\hat{\lambda}_{1},\hat{n})$, ... $\mathbb{I}(\hat{\lambda}_{i},\hat{n})$, ...) corresponding to each of the RoPE dimensions from the $1^{st}$ through the $d^{th}$. With these rescaled $\text{RoPE}$ parameters, the target $\text{LLM}$ can be tuned to attain the lowest possible next token prediction loss, denoted as $\mathcal{L}$ (equivalently, the lowest perplexity), on input sequences $\mathbf{X}$ of length $L'$.

\subsection{Searching the Non-uniform Position Interpolation}

In order to address the issue outlined in Eq.~\ref{eq:problem}, we propose a straightforward but robust approach that identifies the best RoPE rescale factors. This method is specifically designed to take full advantage of the complex, multidimensional variations present in position embedding.

\textbf{Search space}. We construct an extensive search space that accounts for both forms of non-uniformity. Table~\ref{tbl:searchspace} presents a detailed overview of this search space. In particular, we enable independent optimization of the rescale factor for each dimension within the RoPE embedding. To simplify the configuration of the search space, we opt to search directly over $\lambda_i$ and $\hat{n}$ as proxies for $\mathbb{I}(\hat{\lambda}_{i},\hat{n})$, with the relationship $\hat{\lambda}_{i}=1/\lambda_i$. As detailed in Table~\ref{tbl:searchspace}, $\lambda_i$ ranges from a minimum of 1.0 (representing pure extrapolation) up to $s\times1.25$ (corresponding to interpolation greater than that of PI), incremented by 0.01. Here, $s$ denotes the ratio by which the desired context window is extended.

The parameter $\hat{n}$ determines how many of the earliest token positions preserve the original RoPE embedding and are not modified by position interpolation. In practice, we set $\hat{n}$ to explore values from the set \{0, 1, 2, 4, 8, 12, 16, 20, 24, 28, 32, 64, 128, 256\}. Setting $\hat{n}=0$ results in the application of the learned rescale factors to every token position.

\textbf{Evolution-based search}. The search space detailed in Table~\ref{tbl:searchspace} encompasses a vast array of possible positional interpolation configurations, making comprehensive exploration highly challenging. As an illustration, employing an extension ratio of $s=4\times$ results in $400^{128/2}\times14=4\times10^{167}$ potential combinations. The size of the search space increases exponentially with greater extension ratios. To cope with this complexity, we adopt an evolution search approach~\cite{guo2020single} and incorporate two optimization strategies to significantly improve the efficiency of the search process. The overall procedure is outlined in Algorithm~\ref{alg:search}.

\textit{Optimized initial population generation}. Rather than selecting all $P$ rescale factors at random to form the initial population, we explicitly include the three RoPE rescale factors representing PI, NTK, and YaRN as members of this population. The rest of the $P$-3 individuals are generated by introducing random mutations to these three rescale factors, each with a probability of $p$.

\textit{Monotonically non-decreasing constraint}. Upon producing the initial set of candidates, we evaluate the LLM perplexity for every individual by assigning their respective RoPE rescale factors within the target LLM and measuring the perplexity on input $\mathbf{X}$. The most performant $k$ individuals are selected as parents for subsequent evolutionary steps. Nonetheless, the extensive nature of the search space means that random mutation and crossover often generate suboptimal offspring, resulting in many unproductive perplexity evaluations. This inefficiency becomes especially pronounced when $L'$ is large, due to the computational demands of each individual perplexity inference.

To tackle this issue, we enforce a non-decreasing monotonicity condition on the sampled RoPE rescale coefficients such that $\lambda_i\le \lambda_{i+1}$ holds for all $i$. Only those RoPE configurations adhering to this requirement are utilized when assessing perplexity in the LLM, thereby greatly narrowing the search space. In particular, we stipulate that $\lambda_i$ should monotonically increase with respect to the RoPE dimension index (i.e., $i$ = 0,...,63). The rationale for enforcing this dimensional monotonicity arises from NTK theory~\cite{jacot2018neural,tancik2020fourier,ntk}, which indicates that lower dimensions, corresponding to higher frequencies, necessitate reduced interpolation (smaller $\lambda_i$), while higher dimensions with lower frequencies can tolerate increased interpolation (larger $\lambda_i$).

\begin{algorithm}[t]
	\small
	\caption{The search algorithm}
	\textbf{Input:} target LLM, input samples $\mathbf{X}$, population size $P$, mutation size $N_1$, crossover size $N_2$, max iterations $\mathcal{T}$, mutate probability $p$\\

	\begin{algorithmic}[1]
		\label{alg:search}
		\STATE Initialize Top-k as empty set;
		\STATE Generate initial population $\text{P}_0$ using \textit{Initialize\_population\_with\_optimization} ($P$, $\mathbf{X}$, $p$); 
		\FOR{$i$ from $1$ to $\mathcal{T}$}
		\STATE Evaluate $\text{P}_{i-1}$ on $\mathbf{X}$ via \textit{Compute\_perplexity} (LLM, $\text{P}_{i-1}$, $\mathbf{X}$);
		\STATE  Refresh Top-k with \textit{Update\_Topk} (Top-k, $\text{P}_{i-1}$);
		\STATE Create $\text{P}_{mutation}$ using \textit{Mutation\_with\_mono\_constraint} (Top-k, $N_1$, $p$); 
		\STATE Construct $\text{P}_{crossover}$ from \textit{Crossover\_with\_mono\_constraint} (Top-k, $N_2$);
		\STATE Form new population $\text{P}_i$ by combining $\text{P}_{mutation}$, $\text{P}_{crossover}$, and Top-k;
		\ENDFOR
		\STATE Output the member with minimum perplexity in Top-k;
	\end{algorithmic}
\end{algorithm}

\textbf{8$\times$ extension without fine-tuning}. Through evolutionary search, our approach efficiently discovers non-uniform RoPE rescale factors that retain essential dimensions and positions, thereby reducing information loss caused by interpolation. As shown in Fig.\ref{fig:non-ft}, this technique allows us to expand LLaMA2’s context window from 4k to 32k tokens without the need for fine-tuning. By comparison, alternative approaches like PI, as well as non-uniform NTK and YaRN, exhibit a sharp increase in perplexity once the context window is doubled.

\subsection{Extending LLM Context Window to 2048K}
\label{sec:extendto2048k}

\textbf{Progressive extension to 2048k}. Here, we present our strategy for scaling the context window of pre-trained LLMs well beyond the standard 4k, achieving sizes up to and surpassing 2048k. Our proposed non-uniform positional interpolation technique allows for an 8$\times$ context window increase without the need for any additional fine-tuning. To reach even greater extensions, such as 512$\times$, fine-tuning becomes essential. One approach involves identifying suitable RoPE rescaling factors for the desired 2048k window, followed by fine-tuning the model. Nevertheless, this process is hindered by significant computational and resource costs. Furthermore, as highlighted by our empirical observations, effectively fine-tuning LLMs at such high extension scales presents notable difficulties (refer to Appendix).

Fortunately, LongRoPE demonstrates robust performance on both the base model and the extended LLM after fine-tuning. To this end, we propose a resource-efficient, incremental approach that reaches the desired 2048k context length in only 1k fine-tuning iterations, while limiting the maximum training sequence to 256k tokens.

$\Diamond$ \textit{Investigating LongRoPE Search for Scaling Pre-trained LLMs to 256k}. Using LLaMA2 as a reference model, we systematically explore expansions to context window lengths of 128k and 256k tokens. In this phase, the context is enlarged by factors of 32$\times$ and 64$\times$, correspondingly.

$\Diamond$ \textit{Fine-tuning to 256k}. Subsequently, we adapt the pre-trained LLM for a 256k context window via fine-tuning. The process involves an initial fine-tuning of LLaMA2 for 400 steps employing RoPE rescaled factors set for 128k. After these steps, we swap in the RoPE rescaled factors for 256k at the obtained checkpoint and carry out a further 600 fine-tuning steps. Compared to direct fine-tuning at 256k, this strategy demonstrates greater efficiency.

$\Diamond$ \textit{Scaling fine-tuned extended LLM to 2048k via LongRoPE search}. In the final step, we apply an additional search procedure on the LLM previously fine-tuned at a 256k sequence length. This approach achieves an expansive context window of 2048k tokens, all without needing any supplementary fine-tuning. The resulting extension ratio amounts to $512\times$.

\textbf{Performance on shorter context windows}.  Upon increasing the context window to a vast 2048k length, a decline in performance emerges within the original, shorter context window. This phenomenon, previously observed with positional interpolation~\cite{pi}, arises because the technique compacts position embeddings corresponding to the original context window into a considerably smaller region in the high-dimensional embedding space, which in turn impairs the language model's effectiveness. When applying a 512$\times$ extension ratio, this effect is especially pronounced, leading to severe crowding of positions within the initial 4k context window.

To address this issue, an additional evolution search is conducted on the extended LLM, specifically targeting the optimization of RoPE rescale factors at shorter context lengths such as 4k and 8k. Since less positional interpolation is necessary for these lengths, the upper limit for the searched $\lambda$ values is reduced accordingly. The LLM then adapts the respective RoPE rescale factors on-the-fly during inference.

\section{LongRoPE}

\label{sec:method}
Building on these insights, we propose LongRoPE, which initially incorporates an effective search strategy designed to leverage both forms of non-uniformity. Utilizing this approach, we further extend the context window of LLMs to exceed 2 million tokens.

\subsection{Problem Formulation}

These two forms of non-uniformity result in an extensive set of possible solutions and increase the difficulty of the optimization process. To tackle this challenge, we recast the multidimensional non-uniform position interpolation optimization task as a search problem.

For a LLM targeting a  context window size of $L'$ and lengthy input documents $\mathbf{X}$, where each $\mathbf{x}\in\mathbf{X}$ surpasses $L'$ in token length, we denote the original rotary angle of the $i^{th}$ dimension in RoPE embedding at token position $n$ as  $\frac{n}{\beta^i}$. The optimization problem is then formulated as follows:

\begin{equation}
\begin{gathered}
		\label{eq:problem}

\underset {\mathbf{x}\in\mathbf{X}; \, |\mathbf{x}|\ge L'}{ \operatorname {arg\,min} } \,  \mathcal{L}\, \left( \text{LLM} (\text{RoPE}, \mathbf{X})\right),	\text{where \;} 
\\
\scalebox{0.93}{
$\underset {\begin{subarray}{l} i=0,\cdot\cdot,\frac{d}{2}-1;\\ n\in[ 0,|\mathbf{x}|);\end{subarray}}{\text{RoPE}_(n)}= {\left[\cdot\cdot, \cos\left(\mathbb{I}(\hat{\lambda}_{i}, \hat{n})\cdot\frac{n}{\beta^i}\right),  \sin\left(\mathbb{I}(\hat{\lambda}_{i}, \hat{n})\cdot\frac{n}{\beta^i}\right),\cdot\cdot\right] }$}\\
	\text{where \;} \mathbb{I}(\hat{\lambda}_{i},\hat{n})=\begin{cases}
	1&\mbox{\;  $n<\hat{n}$}\\
{\frac{1}{\lambda}_{i}}&\mbox{\; $n\geq\hat{n}$}
\end{cases}
	\end{gathered}
\end{equation}

In this approach, we define a group of rescale factors, $\mathbb{I}(\hat{\lambda}_{i},\hat{n})$, which are designed to address both types of non-uniformities present. Here, $\hat{\lambda_i}$ and $\hat{n}$ represent, respectively, the degree of non-uniformity in the RoPE dimensions and the token positions. In particular, $\mathbb{I}(\hat{\lambda}_{i},\hat{n})$ is utilized to adjust the rotation angle for the $i^{th}$ RoPE dimension, applying the specific rescale parameter $\hat\lambda_{i}$ according to a position threshold $\hat{n}$. For the first $\hat{n}$ token positions, the original rotary angle $\frac{n}{\beta^i}$ from RoPE remains unchanged, meaning the rescale factor $\hat\lambda_{i}$ is not active. Once token positions reach $n\ge\hat{n}$, the rescaling begins to affect the rotary angle computation.

With a desired context window length of $L'$, our task is to determine the best rescale factors ($\mathbb{I}(\hat{\lambda}_{0},\hat{n})$, $\mathbb{I}(\hat{\lambda}_{1},\hat{n})$, ..., $\mathbb{I}(\hat{\lambda}_{i},\hat{n})$, ...) for each RoPE dimension from the $1^{st}$ through $d^{th}$. The ultimate goal is that, by applying these optimal rescalings to the $\text{RoPE}$ within the target $\text{LLM}$, the next token prediction loss, $\mathcal{L}$—that is, the perplexity—is minimized when the model processes input samples $\mathbf{X}$ of length $L'$.

\subsection{Searching the Non-uniform Position Interpolation}

To address the challenge outlined in Eq.~\ref{eq:problem}, we propose a straightforward but powerful technique that identifies the optimal RoPE rescaling coefficients, thereby maximizing the utilization of multidimensional non-uniform characteristics in position embedding.

\textbf{Search space}. We construct an extensive search space that incorporates the two types of non-uniformity under consideration. The search space configuration is detailed in Table~\ref{tbl:searchspace}. In particular, our approach permits individual tuning of the rescale factor for each dimension within the RoPE embedding. To streamline the design of the search space, we directly optimize over $\lambda_i$ and $\hat{n}$, rather than $\mathbb{I}(\hat{\lambda}_{i},\hat{n})$, with $\hat{\lambda}_{i}=1/\lambda_i$. As depicted in Table~\ref{tbl:searchspace}, the permissible range for $\lambda_i$ starts from 1.0 (which represents straightforward extrapolation) and extends to $s\times1.25$ (which allows for broader interpolation than PI), incrementing in steps of 0.01, where $s$ refers to the intended ratio for context window expansion.

The parameter $\hat{n}$ determines how many of the initial token positions maintain the original RoPE embedding, foregoing any positional interpolation. In practice, $\hat{n}$ is selected from the set \{0, 1, 2, 4, 8, 12, 16, 20, 24, 28, 32, 64, 128, 256\}. If $\hat{n}=0$, then all token positions are assigned the optimized rescale factors.

\textbf{Evolution-based search}. The extensive search space detailed in Table~\ref{tbl:searchspace} encompasses a wide array of positional interpolation methods, rendering exhaustive search impractical. For instance, with an extension factor of $s=4\times$, the potential combinations are as high as $400^{128/2}\times14=4\times10^{167}$. This search space grows exponentially with increased extension ratios. To effectively navigate this complexity, we employ evolution search~\cite{guo2020single}, further enhancing search efficiency through two specifically designed optimization strategies. The complete search workflow is depicted in Algorithm~\ref{alg:search}.

\textit{Optimized initial population generation}. Rather than assigning random values to all $P$ rescale factors in the starting population, we explicitly include the three RoPE rescale factors used in PI, NTK, and YaRN as separate members from the outset. The remaining $P-3$ individuals are produced by randomly applying mutations to these three rescale factors, each with probability $p$.

\textit{Monotonically non-decreasing constraint}. Once the initial population is formed, we evaluate each individual by measuring its LLM perplexity. This involves adjusting the target LLM with the assigned RoPE rescale factors and calculating the perplexity on the input $\mathbf{X}$. The highest-performing $k$ individuals are then selected as parents for the evolutionary process. Due to the immense size of the search space, straightforward application of mutation and crossover often results in the exploration of suboptimal solutions, which entails superfluous perplexity evaluations. This inefficiency is exacerbated when $L'$ is large, as each perplexity computation is computationally intensive.

To manage this, a non-decreasing monotonicity restriction is enforced on the sampled RoPE rescaling factors, formally written as $\lambda_i \le \lambda_{i+1}$. Only those RoPE variants adhering to this monotonic property are considered during LLM perplexity measurements, thereby dramatically cutting down the search complexity. More concretely, this setup enforces that $\lambda_i$ should grow monotonically as the RoPE dimension $i$ increases (with $i=0, \ldots, 63$). This dimensional ordering is justified by NTK theory~\cite{jacot2018neural,tancik2020fourier,ntk}, which indicates that lower-indexed dimensions—associated with higher frequencies—warrant less interpolation (implying a smaller $\lambda_i$), while dimensions with higher indices—associated with lower frequencies—permit greater interpolation (and thus, a larger $\lambda_i$).

\begin{algorithm}[t]
	\small
	\caption{The search algorithm}
	\textbf{Input:} target LLM, input samples $\mathbf{X}$, population size $P$, mutation size $N_1$, crossover size $N_2$, maximum number of iterations $\mathcal{T}$, mutation probability $p$\\

	\begin{algorithmic}[1]
		\label{alg:search}
		\STATE Top-k $\leftarrow$ $\phi$;
		\STATE $\text{P}_0$ $\leftarrow$ \textit{Initialize\_population\_with\_optimization} ($P$, $\mathbf{X}$, $p$);
		\FOR{$i$ = $1$ to $\mathcal{T}$}
		\STATE \textit{Compute\_perplexity} (LLM, $\text{P}_{i-1}$, $\mathbf{X}$);
		\STATE  Top-k $\leftarrow$ \textit{Update\_Topk} (Top-k, $\text{P}_{i-1}$);
		\STATE $\text{P}_{mutation}$ $\leftarrow$ \textit{Mutation\_with\_mono\_constraint} (Top-k, $N_1$, $p$);
		\STATE $\text{P}_{crossover}$ $\leftarrow$ \textit{Crossover\_with\_mono\_constraint} (Top-k, $N_2$);
		\STATE $\text{P}_i$ $=$ $\text{P}_{mutation}$ $\cup$ $\text{P}_{crossover}$ $\cup$ Top-k;
		\ENDFOR
		\STATE Output the member with the minimum perplexity in Top-k;
	\end{algorithmic}
\end{algorithm}

\textbf{8$\times$ extension without fine-tuning}. By leveraging evolutionary search, our approach successfully determines non-uniform RoPE rescale factors that retain crucial positions and dimensions, thereby mitigating information loss due to interpolation. As illustrated in Fig.\ref{fig:non-ft}, this enables us to expand the context window of LLaMA2 from 4k to 32k tokens without any fine-tuning. In comparison, other approaches—such as PI, and non-uniform NTK and YaRN—result in a sharp increase in perplexity once the context window exceeds a 2$\times$ extension.

\subsection{Extending LLM Context Window to 2048K}
\label{sec:extendto2048k}

\textbf{Progressive extension to 2048k}. In this section, we present our approach for expanding the context window of pre-trained LLMs from the typical 4k up to beyond 2048k. As shown previously, our non-uniform positional interpolation technique can provide an 8$\times$ increase without the need for fine-tuning. However, scaling to larger windows (such as a 512$\times$ extension) necessitates fine-tuning the models. One possible strategy involves identifying suitable RoPE rescale factors tailored for the intended 2048k window, followed by a fine-tuning process. Nonetheless, this approach encounters significant obstacles due to the substantial computational demands required. Additionally, our observations indicate that effectively fine-tuning LLMs at such high extension ratios is quite challenging (refer to the Appendix for details).

Fortunately, LongRoPE performs well on both the base model and after it has undergone extended fine-tuning. To this end, we propose a resource-efficient, gradual approach that reaches the desired 2048k target using only 1k fine-tuning steps, all confined to a 256k training sequence length.

$\Diamond$ \textit{LongRoPE search for expanding pre-trained LLM to 256k context length}. Using LLaMA2 as a case study, we perform a search targeting context window sizes of 128k and 256k. At this phase, the extension ratios are 32$\times$ and 64$\times$, accordingly.

$\Diamond$ \textit{Fine-tuning to 256k}. Subsequently, the pre-trained LLM is further fine-tuned to support a context window of 256k tokens. To accomplish this, we initially perform 400 fine-tuning steps on LLaMA2 utilizing RoPE rescaling parameters corresponding to 128k. Following this, we update the finished checkpoint by switching the RoPE rescaling factors to those suitable for 256k, and proceed with another 600 fine-tuning steps. This staged approach is observed to be more efficient than immediate fine-tuning to 256k.

$\Diamond$ \textit{Expanding fine-tuned extended LLM to 2048k via LongRoPE search}. In the last step, we apply an additional search on the LLM previously fine-tuned to support 256k sequence lengths. This procedure enables the model to reach a massive context window of 2048k tokens, all without additional fine-tuning. In total, this reflects an extension factor of $512\times$.

\textbf{Restoration within the shorter context window}. Upon extending the context window to an exceptionally large 2048k tokens, we observe a degradation in performance for sequences situated in the initial context range. This phenomenon is well-documented in positional interpolation~\cite{pi}, where embedding positions from the shorter context are constrained into a compressed segment of the representation space when mapped to higher dimensions, thereby impairing the efficacy of the language model. At an extension ratio of 512$\times$, the embeddings of positions inside the standard 4k context window become densely packed, exacerbating this issue.

To address this issue, we carry out an additional evolution-based search on the expanded LLM, fine-tuning the RoPE rescale parameters for shorter context lengths such as 4k and 8k. We decrease the upper bound for the searched $\lambda$ since these shorter lengths necessitate less positional interpolation. At inference time, the LLM adapts the relevant RoPE rescale parameters on-the-fly.

 \section{Experiments}

\label{sec:exp}
\subsection{Setup}
\textbf{Evaluation Tasks and models}. Our experiments incorporate LongRoPE into both LLaMA2-7B and Mistral-7B, assessing model performance from three different perspectives: (1) perplexity when handling lengthy documents to gauge effectiveness in extended contexts; (2) a Passkey retrieval task that tests the model's capability to extract a specific passkey embedded among predominantly irrelevant content; and (3) established LLM benchmark tasks restricted to a short context window of 4096 tokens.

\textbf{Fine-tuning}. In the case of LLaMA2, fine-tuning is performed using a linear-decayed learning rate set at 2e-5, alongside a global batch size of 32. Training proceeds for 400 iterations utilizing the Redpajama~\cite{redpajama} dataset, which is divided into 128k-length segments and bounded by BOS and EOS markers. Subsequently, using the resulting checkpoint, training is extended for an extra 600 steps in order to accommodate a 256k context window. The process for the 128k context size is distributed across 8 A100 GPUs leveraging the system described in~\cite{cube}, whereas scaling to 256k necessitates 16 A100 GPUs. For the Mistral model, fine-tuning adopts a fixed learning rate of 1e-6, with a global batch size of 64. Both the 128k and 256k variants utilize the YaRN~\cite{yarn} protocol, involving 400 training steps on Together Computer's Long-Data Collections~\cite{mistraldata}, each using sequences of 16k tokens, and training is executed on 4 A100 GPUs.

\textbf{Search}. When targeting window sizes up to 256k, we configure the search with $P$=64, $N_1$=$N_2$=16, $p$=0.3, and $\mathcal{T}$=40, and each iteration involves selecting the top-32 candidates for mutation and crossover operations. Perplexity is evaluated based on 5 randomly chosen samples from the PG19 validation set, each at least as long as the intended context window. For window sizes greater than 512k, the population size, as well as the numbers for mutation and crossover candidates, are reduced by half. In this setting, perplexity is computed using 3 random validation samples from Pile-Books3~\cite{pile}.

\textbf{Baselines}. To achieve a context window of 2048k, we conducted fine-tuning on models with either 128k or 256k context lengths, resulting in LongRoPE-2048k (ft=128k) for LLaMA2 and LongRoPE-2048k (ft=256k) for Mistral. Our evaluation involves a comparison among these four models and leading baseline methods for context window extension, with a focus on open-source LLMs that underwent fine-tuning following positional interpolation via techniques such as PI, NTK, and YaRN. Included in our set of comparison models are Together-32k~\cite{together}, Code LLaMA~\cite{codellama}, LongLoRA-full-FT-100k~\cite{longlora}, YaRN-LLaMA, and YaRN-Mistral~\cite{yarn}.

\subsection{Main Results}

\label{sec:mainresult}
\textbf{Long sequence language modeling up to 256k tokens}. Our analysis starts with a comparison to leading long-context LLMs when evaluated on sequences of up to 256k tokens. To illustrate the robustness of our approach, we employ two datasets: the Proof-pile~\cite{pg19} and the PG19~\cite{pile} test sets. Perplexity scores are computed across a range of context sizes, using a sliding window of 256 tokens. For PG19, evaluation includes all 100 test documents. In line with YaRN~\cite{yarn}, we sample 10 random segments—each no shorter than 128k tokens—from the Proof-pile dataset.

Table~\ref{tbl:proof-pile} and Table~\ref{tbl:pg19} present a comparison of perplexity scores for LLaMA2 and Mistral models, each enhanced using various interpolation techniques, evaluated on Proof-pile and PG19 datasets, respectively. Two principal findings emerge: \textbf{(1)} the perplexity of our extended models consistently decreases as evaluation length increases from 4k to 256k, indicating effective exploitation of extended contexts. \textbf{(2)} Despite utilizing a context window up to 16$\times$ longer—which typically poses difficulties for preserving short-context performance—our LongRoPE-2048k models not only maintain but surpass state-of-the-art baselines at context lengths of up to 256k.

\textbf{Long sequence language modeling beyond 2000k}. To assess performance on very large-scale documents, the Books3~\cite{pile} dataset is employed. For practical evaluation, we draw a random sample of 20 books, each with lengths greater than 2048k, and utilize a 256k sliding window for analysis.

As detailed in Table~\ref{tbl:books3}, LongRoPE is able to expand the context windows of both LLaMA2-7B and Mistral-7B to 2048k, while maintaining perplexity on par with or better than baseline models within the 8k-128k range. In comparing the two models at 2048k, clear differences emerge. Mistral demonstrates superior results relative to baselines for shorter contexts, but its perplexity rises above 7 when exceeding 256k. LLaMA2 exhibits expected trends, with perplexity gradually lowering as context length increases, and only slight upticks observed at 1024k and 2048k. For LLaMA2, employing LongRoPE-2048k with a fine-tuning length of 256k provides advantages over a 128k window, attributable to the reduced secondary extension ratio (8$\times$ as opposed to 16$\times$). On the other hand, Mistral yields better outcomes with a 128k fine-tuning window. This discrepancy largely stems from the use of a 16k training length (as recommended by YaRN's configuration) during fine-tuning for both 128k and 256k windows in Mistral, which in turn constrains its ability to increase the context window following fine-tuning.

\textbf{Passkey retrieval}. In this section, we examine how the size of the context window impacts generative tasks. Our approach adopts the synthetic passkey retrieval benchmark introduced by~\cite{passkey}. During each trial, the model must locate a concealed random passkey—specifically, a five-digit number—within an extended document. Full details of the prompt structure are provided in the appendix. For evaluation, we conduct 10 runs of the task, ensuring the passkey is embedded at a randomly selected position that is uniformly sampled throughout the context window.

Figure~\ref{fig:passkey} presents a comparison of retrieval accuracy across different models and baselines. While previous models experience a sharp decline in accuracy—falling to zero once the context exceeds 128k tokens—our LongRoPE-LLaMA2-2048k (ft=256k) consistently achieves high retrieval accuracy ($\ge$90\%) across a wide context range, from 4k up to 2048k tokens, even under the difficult scenario of extracting a passkey from millions of tokens. Similarly, LongRoPE-Mistral-2048k (ft=128k) maintains perfect (100\%) accuracy up to 1800k tokens and then decreases to 60\% at 2048k. This performance trend is consistent with Table~\ref{tbl:books3}, which shows a modest perplexity increase at 2048k.

\textbf{Evaluation on standard benchmarks within the native context window}. To assess the performance of LongRoPE-2048k models within their original context window, we utilize the Hugging Face Open LLM Leaderboard~\cite{huggingfaceboard} for both zero-shot and few-shot tasks. Specifically, evaluations are conducted with 25-shot ARC-Challenge~\cite{allenai:arc}, 10-shot HellaSwag~\cite{zellers2019hellaswag}, 5-shot MMLU~\cite{mmlu}, and 0-shot TruthfulQA~\cite{lin2021truthfulqa}.

As presented in Table~\ref{tbl:commonsense}, our models perform on par with the initial benchmark intended for shorter context windows, and even surpass the original Mistral on TruthfulQA by an increase of +0.5\%. While LongRoPE-LLaMA2-2048k, which underwent fine-tuning at 256k, experiences marginally higher declines in performance, its results for the majority of tasks are still largely acceptable.

\subsection{Ablation Results}

\textbf{Effectiveness of the second positional interpolation}. Within the framework of our progressive extension methodology, we apply our search algorithm to perform a second stage of non-uniform positional interpolation on extended and fine-tuned LLMs. To assess the impact of this approach, we execute experiments using our fine-tuned LLaMA2-256k model, expanding its context length to 512k, 1024k, and 2048k via both PI and YaRN. According to the data presented in Table~\ref{tbl:secondsearch}, our approach using non-uniform positional interpolation maintains stable perplexity scores. By contrast, models employing PI and YaRN exhibit a rapid increase in perplexity as the extension factor grows.

\textbf{Recovery performance at reduced context lengths.} To address the drop in effectiveness observed at lower context lengths, we recalibrate the RoPE scaling parameters in LongRoPE-2048k using our optimization procedure. In particular, by restricting the upper bounds of the scaling factors considered in the search, we promote reduced levels of interpolation for context windows of 4k and 8k. Table~\ref{tbl:shortperformance} presents perplexity outcomes for LongRoPE-LLaMA2-2048k on the Proof-pile at 4k and 8k contexts, as well as average scores across LLM benchmarks. These findings clearly indicate notable gains in performance when operating on shorter context lengths.

\textbf{Analysis on the two forms of non-uniformities}. To assess the individual effects of each non-uniformity, we conduct an ablation study to determine their respective contributions to model performance. Our experimental setup consists of two parts: (i) expanding LLaMA2-7B to shorter context lengths (16k and 32k) via several strategies—Position Interpolation (PI), searching exclusively over the RoPE dimension, and searching with respect to both non-uniformities; and (ii) extending our 256k-length fine-tuned LLaMA2 to 2048k using the same methods. All perplexity measurements are taken without additional fine-tuning. According to Table~\ref{tbl:searchdimension}, introducing non-uniformity into the RoPE dimension leads to a notable reduction in perplexity when compared to the standard linear interpolation used by PI. At 16k and 32k lengths, the non-uniform token position contributes clear performance gains, though this benefit does not extend to 2048k—possibly due to challenges inherent in such long sequences. Maintaining the initial tokens without applying interpolation offers little value at this scale, indicating an avenue for further investigation.

\section{Related Works}

Besides approaches centered on position interpolation, this section examines alternative methodologies from the literature.

\textbf{Retrieval-based} approaches employ external memory systems to store extensive historical context and utilize retrieval mechanisms to fetch pertinent documents during inference~\cite{longllama,wang2023augmenting,borgeaud2022improving}. Implementing these strategies generally requires substantial architectural adjustments to LLMs. By comparison, our method introduces only minimal changes, specifically in the positional embedding layer, resulting in a more streamlined solution. Furthermore, our approach accommodates a broader set of long-context tasks beyond just retrieval, encompassing applications such as summarizing lengthy documents and enabling few-shot learning.

\textbf{Attention-based context window extensions}. In addition to positional embedding interpolation, alternative studies have managed to extend the input context for the original LLM context window by altering attention mechanisms~\cite{han2023lminfinite, streamingllm,parallelwindow}. Central to these methods is addressing the problem of escalating attention complexity from introducing new positions through innovative attention masks. These strategies can be used alongside positional interpolation, as the two approaches are complementary.

\textbf{Fine-tuning based approaches} investigate strategies for adapting pre-trained LLMs to support extended context lengths by altering position embeddings. Approaches such as Code LLaMA~\cite{codellama}, LLaMA2 Long~\cite{xiong2023effective}, and ScaledRoPE~\cite{liu2023scaling} employ a significantly increased RoPE base parameter and subsequently fine-tune on the desired longer sequence lengths. In contrast, our approach accommodates a diverse range of target lengths and is capable of reaching beyond 2M tokens. With the increasing memory requirements for fine-tuning at extended context lengths (exceeding 128k), solutions like LongLoRA~\cite{longlora} and PoSE~\cite{zhu2023pose} have been introduced to address the corresponding GPU overhead. Our method is complementary to such efficient fine-tuning techniques.

\section{Conclusion}

In this study, we introduce LongRoPE, a technique that significantly increases the context window of LLMs up to 2048k tokens, all the while preserving their performance in the shorter, original context range. Our approach takes advantage of two distinct non-uniform patterns present in the RoPE positional embeddings, identified via a streamlined evolutionary search procedure. This yields dual advantages: it offers effective initialization for subsequent fine-tuning and supports an 8$\times$ enlargement of the context window even in the absence of additional fine-tuning. Furthermore, we outline a progressive extension methodology, whereby LLMs fine-tuned with 256k-length contexts are used as a foundation to achieve a 2048k-token context window without needing further fine-tuning steps. Comprehensive empirical evaluations affirm the efficacy of LongRoPE. We anticipate that the LongRoPE-2048k models will unlock a range of novel applications requiring extended context and stimulate ongoing research in this direction.

\section*{Broader Impacts} The research described in this paper aims to further the development of the Machine Learning discipline. While our contributions may influence society in various ways, we do not identify any particular issues that require explicit attention at this time.

	\balance

\end{document}